% Source: https://artofproblemsolving.com/community/c3h1458123_triangle_inscribed_circle
% Stable identity: aops-topic-1458123-q01
Let $a=BC$, $b=CA$, $c=AB$ and $s=\dfrac{a+b+c}{2}$. Write $x=AQ=AR$, $y=BR=BP$, $z=CP=CQ$. Since $R$ is on $AB$, $P$ on $BC$ and $Q$ on $CA$,
\[
x+y=c,\qquad y+z=a,\qquad z+x=b .
\]
Adding, $x+y+z=s$, so
\[
x=s-a,\qquad y=s-b,\qquad z=s-c .
\]
Now let the incircle touch $BC$, $CA$, $AB$ at $P'$, $Q'$, $R'$. The two tangent segments from a point to a circle are equal, so $AQ'=AR'$, $BR'=BP'$, $CP'=CQ'$. The same computation shows $AR'=s-a$, $BP'=s-b$, $CQ'=s-c$.

So $R$ and $R'$ are both the point of side $AB$ at distance $s-a$ from $A$, that is, $R=R'$. In the same way $P=P'$ and $Q=Q'$. Thus $P$, $Q$, $R$ are the points where the incircle touches the sides; in particular the incircle passes through all three. $\blacksquare$

\begin{asy}
usepackage("amsmath");
size(6cm);
defaultpen(fontsize(10pt));
pair ft(pair P,pair U,pair V){return U+dot(P-U,V-U)/dot(V-U,V-U)*(V-U);}
path ra(pair P,pair u,pair v,real s=0.2){return P+s*unit(u)--P+s*unit(u)+s*unit(v)--P+s*unit(v);}
pair A=(1.3,4.2),B=(0,0),C=(6,0);
real a=length(B-C),bb=length(C-A),c=length(A-B),s=(a+bb+c)/2;
pair I=(a*A+bb*B+c*C)/(a+bb+c);
pair P=B+(s-bb)*unit(C-B),Q=C+(s-c)*unit(A-C),R=A+(s-a)*unit(B-A);
real r=length(P-I);
draw(circle(I,r),blue);
draw(I--P,gray+dashed); draw(I--Q,gray+dashed); draw(I--R,gray+dashed);
draw(A--B--C--cycle,linewidth(0.9));
dot(A);dot(B);dot(C);dot(P);dot(Q);dot(R);dot(I,blue);
label("$A$",A,N); label("$B$",B,SW); label("$C$",C,SE);
label("$P$",P,S); label("$Q$",Q,NE); label("$R$",R,W); label("$I$",I,NE,blue);
label("$x$",(A+R)/2,W,red); label("$x$",(A+Q)/2,NE,red);
label("$y$",(B+R)/2,W,red); label("$y$",(B+P)/2,S,red);
label("$z$",(C+P)/2,S,red); label("$z$",(C+Q)/2,NE,red);
label("\mbox{}",(min(currentpicture,user=true)+max(currentpicture,user=true))/2,black); // keeps the text after the figure black when LaTeX typesets the labels inline
\end{asy}
